\section{Theoretical Setup}

\subsection{Anchored common-chain presentations}

Fix integers $N\geq 1$ and $q\geq 0$, a center $c\in\mathbb R$, and a
half-width $h>0$.  Write
\[
\Theta=[c-h,c+h],
\qquad
x=x(\theta):=\frac{\theta-c}{h}\in[-1,1].
\]
The ambient parameter dimension is $p=1$.  For a polynomial $S$, let
$\|\operatorname{coeff}(S)\|_1$ denote the sum of the absolute values of its
coefficients in the standard monomial basis.

Let $\eta=(\eta_1,\ldots,\eta_q)$ be a common triangular chain on
$[-1,1]$.  For $1\leq j\leq q$, let $P_j$ be its chain polynomial in
$(x,y_1,\ldots,y_j)$.  For $1\leq i\leq N$, let $Q_i$ be an output
polynomial in $(x,y_1,\ldots,y_q)$.  We use the one-variable convention in
\citep{balcan2025pfaffian}:
\[
M:=\max_{1\leq j\leq q}\deg P_j\quad(q\geq1),
\qquad
\Delta:=\max_{1\leq i\leq N}\deg Q_i,
\]
where degrees are total degrees, and $M=0$ when $q=0$.  Define
\[
G_i(x):=Q_i(x,\eta_1(x),\ldots,\eta_q(x)),
\qquad
F_i(\theta):=G_i(x(\theta)),
\]
and
\[
B_P:=
\begin{cases}
\displaystyle\max_{1\leq j\leq q}
 \|\operatorname{coeff}(P_j)\|_1,&q\geq1,\\
0,&q=0,
\end{cases}
\qquad
B_Q:=\max_{1\leq i\leq N}
 \|\operatorname{coeff}(Q_i)\|_1.
\]
The affine change from $x$ to $\theta$ leaves $q,M,\Delta$ unchanged and
satisfies $d/d\theta=h^{-1}d/dx$.

For $R>0$ and $0<\kappa<\infty$, let $\mathcal D_{N,R,\kappa}$ be the class
of Borel probability laws on $\mathbb R^N$ with a Lebesgue density supported
on $[-R,R]^N$ and bounded almost everywhere by $\kappa$.  No product
structure is part of this definition.

\subsection{Assumptions}

\begin{assumption}[Primitive parameter regime]
\label{assump:parameter-regime}
$N\geq1$, $q\geq0$, $h>0$, $R>0$, and $0<\kappa<\infty$;
$\Theta=[c-h,c+h]$; the class $\mathcal D_{N,R,\kappa}$ is nonempty; and all
polynomial degrees and coefficient budgets displayed above are finite static
presentation data.
\end{assumption}

\begin{assumption}[Balcan common-chain presentation]
\label{assump:balcan-common-chain}
Each $\eta_j\in C^1([-1,1])$ satisfies
\[
\eta_j'(x)=P_j(x,\eta_1(x),\ldots,\eta_j(x)),
\qquad \deg P_j\leq M,
\]
and every output has the common-chain representation
\[
G_i(x)=Q_i(x,\eta_1(x),\ldots,\eta_q(x)),
\qquad \deg Q_i\leq\Delta.
\]
Thus $q$ is chain length, $M$ is chain degree, $\Delta$ is output degree,
and $p=1$.
\end{assumption}

\begin{assumption}[Literal anchor and unit-range certificate]
\label{assump:anchored-unit-range}
For every $x\in[-1,1]$, $|\eta_j(x)|\leq1$ for all $j$, and
$Q_1\equiv1$.  No separate lower norm margin or projective-speed bound is
assumed.
\end{assumption}

\begin{assumption}[Arbitrarily correlated capped joint laws]
\label{assump:cube-density-laws}
In every probabilistic clause, $\alpha$ is drawn from an arbitrary
$\mu\in\mathcal D_{N,R,\kappa}$.  Only cube support and the cap on the full
joint density are available; the coordinates may be arbitrarily correlated.
\end{assumption}

\begin{assumption}[Deterministic affine offset and pivot cover]
\label{assump:affine-chart-data}
Whenever the affine clause is invoked, $F_0\in C^1(\Theta)$, and the selected
measurable sets $E_1,\ldots,E_N$ form a partition
$I=\bigsqcup_{j=1}^N E_j$ of the interval under consideration, with
$F_j(\theta)\neq0$ on $E_j$.
\end{assumption}

\subsection{Objective}

Our goal is a single theorem that derives projective conditioning from the
anchored presentation, controls central and affine coefficient incidences for
arbitrary capped joint laws, recovers the exact monic-polynomial baseline with
the leading coefficient kept deterministic, and records the sharp metric
scale forced by the two-feature example.  The result applies to the declared
anchored, unit-range, coefficient-controlled presentation; it makes no claim
that every raw Pfaffian presentation admits such a normalization with
polynomial budgets.
